\documentclass{article}
\usepackage{amsmath,amssymb,amsthm}
\usepackage{algorithm}
\usepackage{algorithmic}
\usepackage{hyperref}
\usepackage{enumitem}
\usepackage{bm}
\usepackage{geometry}
\geometry{margin=1in}
\newtheorem{theorem}{Theorem}
\newtheorem{lemma}{Lemma}
\newtheorem{assumption}{Assumption}
\newtheorem{corollary}{Corollary}
\newcommand{\E}{\mathbb{E}}
\newcommand{\R}{\mathbb{R}}

\begin{document}

\section*{FedProx Local SGD}

\section*{Mathematical Formulation}
Let there be $K$ participating clients.  
Each client $k$ possesses a private convex loss $f_k:\R^d\to\R$ that is $L$‑smooth.  
The global objective is 
\[
F(x)\;:=\;\frac1K\sum_{k=1}^{K}f_k(x),\qquad x\in\R^{d}.
\]

For communication round $t=0,1,\dots,T-1$ the server holds a global model $x_t$.  
Each client receives $x_t$ and performs $E$ local stochastic gradient steps on the
proximal objective
\[
\phi_k^{(t)}(x)\;:=\;f_k(x)+\frac{\mu}{2}\|x-x_t\|^{2},\qquad \mu\ge0 .
\]

Denote by $g_{t}^{k,e}$ an unbiased stochastic gradient of $f_k$ at the local iterate
$x_{t}^{k,e}$, i.e.
\[
\E\!\left[g_{t}^{k,e}\mid\mathcal{F}_{t}^{k,e}\right]=\nabla f_k(x_{t}^{k,e}),
\]
where $\mathcal{F}_{t}^{k,e}$ is the $\sigma$‑algebra generated by all randomness up to
step $e$ of client $k$ in round $t$.
The local update rule is
\begin{align}
x_{t}^{k,0}&=x_t,\\[2mm]
x_{t}^{k,e+1}&=x_{t}^{k,e}
-\eta\Bigl(g_{t}^{k,e}+\mu\bigl(x_{t}^{k,e}-x_t\bigr)\Bigr),
\qquad e=0,\dots,E-1. \label{eq:local_update}
\end{align}
After $E$ steps client $k$ sends $x_{t}^{k,E}$ to the server.  
The server aggregates by simple averaging:
\[
x_{t+1}\;:=\;\frac1K\sum_{k=1}^{K}x_{t}^{k,E}. \tag{Server aggregation}
\]

\section*{Pseudocode}
\begin{algorithm}[H]
\caption{FedProx Local SGD}
\begin{algorithmic}[1]
\REQUIRE Number of clients $K$, local epochs $E$, learning rate $\eta$, proximal weight $\mu$,
initial model $x_0$.
\FOR{communication round $t=0$ {\bf to} $T-1$}
    \STATE Server broadcasts $x_t$ to all clients.
    \FOR{each client $k=1,\dots,K$ \textbf{in parallel}}
        \STATE $x_{t}^{k,0}\leftarrow x_t$.
        \FOR{$e=0$ \textbf{to} $E-1$}
            \STATE Sample mini‑batch $\xi_{t}^{k,e}$ and compute stochastic gradient 
            $g_{t}^{k,e}=\nabla f_k(x_{t}^{k,e};\xi_{t}^{k,e})$.
            \STATE $x_{t}^{k,e+1}\leftarrow x_{t}^{k,e}
                -\eta\bigl(g_{t}^{k,e}+\mu(x_{t}^{k,e}-x_t)\bigr)$.
        \ENDFOR
        \STATE Send $x_{t}^{k,E}$ to the server.
    \ENDFOR
    \STATE $x_{t+1}\leftarrow\frac1K\sum_{k=1}^{K}x_{t}^{k,E}$.
\ENDFOR
\ENSURE Final model $x_T$.
\end{algorithmic}
\end{algorithm}

\section*{Dry Run Example}
Consider a single client ($K=1$) with scalar loss $f(w)=(w-3)^{2}$.  
Thus $\nabla f(w)=2(w-3)$.  
Set $w_0=0$, learning rate $\eta=0.1$, proximal weight $\mu=0$, and $E=3$ local steps.

\begin{center}
\begin{tabular}{c|c|c|c}
Step $e$ & $w^{e}$ & Stochastic gradient $g^{e}$ (here exact) & $w^{e+1}=w^{e}-\eta(g^{e}+\mu(w^{e}-w_0))$\\ \hline
0 & $0$ & $2(0-3)=-6$ & $0-0.1(-6)=0.6$\\
1 & $0.6$ & $2(0.6-3)=-4.8$ & $0.6-0.1(-4.8)=1.08$\\
2 & $1.08$ & $2(1.08-3)=-3.84$ & $1.08-0.1(-3.84)=1.464$\\
\end{tabular}
\end{center}
Objective values: $f(0)=9$, $f(0.6)=5.76$, $f(1.08)=3.6864$, $f(1.464)=2.3410$.  
The iterate moves toward the optimum $w^{\star}=3$ as expected.

\newpage
\section*{Assumptions}
\begin{assumption}[Convexity and Smoothness]\label{ass:convex_smooth}
Each local loss $f_k$ is convex and $L$‑smooth:
\[
\forall x,y\in\R^{d},\qquad
f_k(y)\le f_k(x)+\langle\nabla f_k(x),y-x\rangle+\frac{L}{2}\|y-x\|^{2}.
\]
\end{assumption}

\begin{assumption}[Unbiased Stochastic Gradients]\label{ass:unbiased}
For every client $k$, round $t$, local step $e$,
\[
\E\!\bigl[g_{t}^{k,e}\mid\mathcal{F}_{t}^{k,e}\bigr]=\nabla f_k(x_{t}^{k,e}),
\]
and the variance is uniformly bounded,
\[
\E\!\bigl\|g_{t}^{k,e}-\nabla f_k(x_{t}^{k,e})\bigr\|^{2}\le\sigma^{2}.
\]
\end{assumption}

\begin{assumption}[Learning Rate]\label{ass:lr}
The stepsize satisfies
\[
0<\eta\le\frac{1}{L+\mu}.
\]
\end{assumption}

\section*{Main Convergence Theorem}
\begin{theorem}[Convergence of FedProx Local SGD]\label{thm:main}
Under Assumptions~\ref{ass:convex_smooth}--\ref{ass:lr},
let $x^{\star}\in\arg\min_{x}F(x)$.
After $T$ communication rounds the expected suboptimality obeys
\[
\boxed{\;
\E\!\bigl[F(x_T)-F(x^{\star})\bigr]
\;\le\;
\frac{\|x_0-x^{\star}\|^{2}}{2\eta T}
\;+\;
\frac{\eta L\sigma^{2}}{K}\;+\;
\frac{\mu\|x_0-x^{\star}\|^{2}}{2T}\; }.
\]
In particular, choosing $\eta=\Theta\!\bigl(\frac{1}{\sqrt{T}}\bigr)$ yields
\[
\E\!\bigl[F(x_T)-F(x^{\star})\bigr]=
\mathcal{O}\!\Bigl(\frac{1}{\sqrt{T}}\Bigr)
\quad\text{when }\mu=0,
\]
and $\mathcal{O}(1/T)$ when $\mu>0$ (strongly‑convex–like regime).
\end{theorem}

\section*{Theoretical Properties}
\begin{itemize}
\item \textbf{Stability.} The proximal term $\frac{\mu}{2}\|x-x_t\|^{2}$ enforces
a contractive mapping; larger $\mu$ improves stability at the expense of slower
movement of the global model.
\item \textbf{Hyper‑parameter Sensitivity.} The bound shows a linear dependence on
$\eta$ for the variance term and an inverse dependence on $\eta$ for the
optimization term. The optimal $\eta$ balances these two contributions.
\item \textbf{Comparison to Baselines.}
\begin{itemize}
\item With $\mu=0$ and $E=1$, the algorithm reduces to standard distributed SGD,
and Theorem~\ref{thm:main} recovers the classical $O(1/\sqrt{T})$ rate.
\item For $E>1$ the bound still holds because the proximal term controls drift
between local and global models; this is the essential novelty of FedProx.
\end{itemize}
\end{itemize}

\section*{Discussion}
The theorem demonstrates that FedProx retains the optimal $O(1/\sqrt{T})$ rate
for convex smooth problems while allowing multiple local updates ($E\ge1$).
When the proximal weight $\mu$ is positive the algorithm enjoys a
$\mathcal{O}(1/T)$ rate reminiscent of strongly convex optimisation, even though the
original losses may be only convex. This effect stems from the added quadratic term,
which imparts artificial strong convexity around the current global iterate.

\newpage
\section*{Preliminaries (Definitions and Lemmas)}
\begin{lemma}[One‑step Descent for a Single Client]\label{lem:one_step}
For any client $k$, round $t$, and local step $e$,
\[
\E\!\bigl\|x_{t}^{k,e+1}-x^{\star}\bigr\|^{2}
\le
\bigl(1-2\eta\mu\bigr)\E\!\bigl\|x_{t}^{k,e}-x^{\star}\bigr\|^{2}
-2\eta\E\!\bigl[f_k(x_{t}^{k,e})-f_k(x^{\star})\bigr]
+\eta^{2}\sigma^{2}.
\]
\end{lemma}
\begin{proof}
Starting from the update \eqref{eq:local_update} and expanding the squared norm,
\begin{align}
\|x_{t}^{k,e+1}-x^{\star}\|^{2}
&=\|x_{t}^{k,e}-\eta(g_{t}^{k,e}+\mu(x_{t}^{k,e}-x_t))-x^{\star}\|^{2}\nonumber\\
&=\|x_{t}^{k,e}-x^{\star}\|^{2}
-2\eta\langle g_{t}^{k,e}+\mu(x_{t}^{k,e}-x_t),\,x_{t}^{k,e}-x^{\star}\rangle
+\eta^{2}\|g_{t}^{k,e}+\mu(x_{t}^{k,e}-x_t)\|^{2}. \tag{1}
\end{align}
Taking conditional expectation w.r.t.\ $\mathcal{F}_{t}^{k,e}$ and using unbiasedness
\eqref{ass:unbiased},
\[
\E\!\bigl[g_{t}^{k,e}\mid\mathcal{F}_{t}^{k,e}\bigr]=\nabla f_k(x_{t}^{k,e}).
\]
Hence
\begin{align}
\E\!\bigl[\langle g_{t}^{k,e},x_{t}^{k,e}-x^{\star}\rangle\mid\mathcal{F}_{t}^{k,e}\bigr]
&=\langle\nabla f_k(x_{t}^{k,e}),x_{t}^{k,e}-x^{\star}\rangle. \tag{2}
\end{align}
Using convexity of $f_k$,
\[
f_k(x_{t}^{k,e})-f_k(x^{\star})\le
\langle\nabla f_k(x_{t}^{k,e}),x_{t}^{k,e}-x^{\star}\rangle. \tag{3}
\]
Insert (2) and (3) into (1) and then take full expectation:
\begin{align}
\E\!\bigl\|x_{t}^{k,e+1}-x^{\star}\|^{2}
&\le
\E\!\|x_{t}^{k,e}-x^{\star}\|^{2}
-2\eta\E\!\bigl[f_k(x_{t}^{k,e})-f_k(x^{\star})\bigr] \nonumber\\
&\quad-2\eta\mu\E\!\langle x_{t}^{k,e}-x_t,\,x_{t}^{k,e}-x^{\star}\rangle
+\eta^{2}\E\!\|g_{t}^{k,e}+\mu(x_{t}^{k,e}-x_t)\|^{2}. \tag{4}
\end{align}
Expand the last norm and use $\|a+b\|^{2}\le2\|a\|^{2}+2\|b\|^{2}$:
\begin{align}
\|g_{t}^{k,e}+\mu(x_{t}^{k,e}-x_t)\|^{2}
\le2\|g_{t}^{k,e}\|^{2}+2\mu^{2}\|x_{t}^{k,e}-x_t\|^{2}. \tag{5}
\end{align}
Taking expectation and applying the variance bound from Assumption~\ref{ass:unbiased},
\[
\E\!\|g_{t}^{k,e}\|^{2}
=\E\!\bigl\|\nabla f_k(x_{t}^{k,e})\bigr\|^{2}
+\E\!\bigl\|g_{t}^{k,e}-\nabla f_k(x_{t}^{k,e})\bigr\|^{2}
\le \E\!\bigl\|\nabla f_k(x_{t}^{k,e})\bigr\|^{2}+\sigma^{2}. \tag{6}
\]
Since $f_k$ is $L$‑smooth, $\|\nabla f_k(x)\|^{2}\le2L\bigl(f_k(x)-f_k(x^{\star})\bigr)$.
Using $0<\eta\le1/(L+\mu)$ and discarding the non‑negative term
$2\mu^{2}\|x_{t}^{k,e}-x_t\|^{2}$ yields
\[
\eta^{2}\E\!\|g_{t}^{k,e}+\mu(x_{t}^{k,e}-x_t)\|^{2}
\le\eta^{2}\sigma^{2}+2\eta^{2}L\E\!\bigl[f_k(x_{t}^{k,e})-f_k(x^{\star})\bigr]. \tag{7}
\]
Insert (7) into (4) and regroup terms:
\[
\E\!\|x_{t}^{k,e+1}-x^{\star}\|^{2}
\le\bigl(1-2\eta\mu\bigr)\E\!\|x_{t}^{k,e}-x^{\star}\|^{2}
-2\eta\bigl(1-L\eta\bigr)\E\!\bigl[f_k(x_{t}^{k,e})-f_k(x^{\star})\bigr]
+\eta^{2}\sigma^{2}. \tag{8}
\]
Because $\eta\le1/(L+\mu)\le1/L$, the factor $(1-L\eta)$ is non‑negative and can be dropped to obtain the claimed inequality. \hfill\qedsymbol
\end{proof}

\section*{Main Proof (Step‑by‑Step Derivation)}
We now prove Theorem~\ref{thm:main}.  The proof is organized into numbered
steps for easy reference.

\begin{proof}
\textbf{[Step 1]} \emph{Apply Lemma~\ref{lem:one_step} recursively over the $E$ local steps.}  
Define $\Delta_{t}^{k,e}:=\E\!\|x_{t}^{k,e}-x^{\star}\|^{2}$.  Lemma~\ref{lem:one_step} gives
\[
\Delta_{t}^{k,e+1}
\le (1-2\eta\mu)\Delta_{t}^{k,e}
-2\eta\E\!\bigl[f_k(x_{t}^{k,e})-f_k(x^{\star})\bigr]
+\eta^{2}\sigma^{2}. \tag{9}
\]
Iterating (9) for $e=0,\dots,E-1$ and telescoping yields
\[
\Delta_{t}^{k,E}
\le (1-2\eta\mu)^{E}\Delta_{t}^{k,0}
-2\eta\sum_{e=0}^{E-1}(1-2\eta\mu)^{E-1-e}
\E\!\bigl[f_k(x_{t}^{k,e})-f_k(x^{\star})\bigr]
+E\eta^{2}\sigma^{2}. \tag{10}
\]

\textbf{[Step 2]} \emph{Upper bound the weighted sum of function values.}  
Since $(1-2\eta\mu)^{E-1-e}\le1$, we drop the weights:
\[
\sum_{e=0}^{E-1}(1-2\eta\mu)^{E-1-e}
\E\!\bigl[f_k(x_{t}^{k,e})-f_k(x^{\star})\bigr]
\ge\sum_{e=0}^{E-1}
\E\!\bigl[f_k(x_{t}^{k,e})-f_k(x^{\star})\bigr]. \tag{11}
\]
Thus (10) simplifies to
\[
\Delta_{t}^{k,E}
\le (1-2\eta\mu)^{E}\Delta_{t}^{k,0}
-2\eta\sum_{e=0}^{E-1}
\E\!\bigl[f_k(x_{t}^{k,e})-f_k(x^{\star})\bigr]
+E\eta^{2}\sigma^{2}. \tag{12}
\]

\textbf{[Step 3]} \emph{Relate the client’s initial distance to the global distance.}  
Because $x_{t}^{k,0}=x_t$, we have $\Delta_{t}^{k,0}= \E\!\|x_t-x^{\star}\|^{2}$, which is common for all $k$.
Averaging (12) over $k=1,\dots,K$ and using the server aggregation $x_{t+1}=K^{-1}\sum_{k}x_{t}^{k,E}$,
\begin{align}
\E\!\|x_{t+1}-x^{\star}\|^{2}
&=\E\!\Bigl\|\frac1K\sum_{k}x_{t}^{k,E}-x^{\star}\Bigr\|^{2}
\le\frac1K\sum_{k}\E\!\|x_{t}^{k,E}-x^{\star}\|^{2} \nonumber\\
&\le (1-2\eta\mu)^{E}\E\!\|x_t-x^{\star}\|^{2}
-2\eta\frac1K\sum_{k}\sum_{e=0}^{E-1}
\E\!\bigl[f_k(x_{t}^{k,e})-f_k(x^{\star})\bigr]
+ \eta^{2}\sigma^{2}E. \tag{13}
\end{align}
The inequality follows from Jensen’s inequality applied to the convex square norm.

\textbf{[Step 4]} \emph{Replace the inner sum of local losses by the global loss.}  
Define the averaged local iterate at step $e$ of round $t$:
\[
\bar{x}_{t}^{e}:=\frac1K\sum_{k}x_{t}^{k,e}.
\]
Convexity of $F$ together with Jensen’s inequality gives
\[
\frac1K\sum_{k}f_k(x_{t}^{k,e})\ge F(\bar{x}_{t}^{e}). \tag{14}
\]
Moreover, because $F$ is convex,
\[
F(\bar{x}_{t}^{e})-F(x^{\star})\ge \frac1K\sum_{k}\bigl[f_k(x_{t}^{k,e})-f_k(x^{\star})\bigr]. \tag{15}
\]
Insert (15) into (13):
\[
\E\!\|x_{t+1}-x^{\star}\|^{2}
\le (1-2\eta\mu)^{E}\E\!\|x_t-x^{\star}\|^{2}
-2\eta\sum_{e=0}^{E-1}\E\!\bigl[F(\bar{x}_{t}^{e})-F(x^{\star})\bigr]
+ \eta^{2}\sigma^{2}E. \tag{16}
\]

\textbf{[Step 5]} \emph{Upper‑bound the contraction factor.}  
Since $0<\eta\le\frac1{L+\mu}$ we have $0<2\eta\mu\le1$, hence
\[
(1-2\eta\mu)^{E}\le 1-2\eta\mu E. \tag{17}
\]
Using (17) in (16) and discarding the non‑negative sum over $e$ (which only improves the bound) yields
\[
\E\!\|x_{t+1}-x^{\star}\|^{2}
\le \bigl(1-2\eta\mu E\bigr)\E\!\|x_t-x^{\star}\|^{2}
+ \eta^{2}\sigma^{2}E. \tag{18}
\]

\textbf{[Step 6]} \emph{Recurrence solution.}  
Unfold (18) for $t=0,\dots,T-1$:
\[
\E\!\|x_{T}-x^{\star}\|^{2}
\le (1-2\eta\mu E)^{T}\|x_0-x^{\star}\|^{2}
+ \eta^{2}\sigma^{2}E\sum_{t=0}^{T-1}(1-2\eta\mu E)^{t}. \tag{19}
\]
The geometric series evaluates to
\[
\sum_{t=0}^{T-1}(1-2\eta\mu E)^{t}
=\frac{1-(1-2\eta\mu E)^{T}}{2\eta\mu E}\le\frac1{2\eta\mu E}. \tag{20}
\]
Plugging (20) into (19) gives
\[
\E\!\|x_{T}-x^{\star}\|^{2}
\le (1-2\eta\mu E)^{T}\|x_0-x^{\star}\|^{2}
+\frac{\eta\sigma^{2}}{2\mu}. \tag{21}
\]

\textbf{[Step 7]} \emph{From distance to function suboptimality.}  
By $L$‑smoothness,
\[
F(x_T)-F(x^{\star})\le\frac{L}{2}\|x_T-x^{\star}\|^{2}. \tag{22}
\]
Taking expectation and using (21):
\[
\E\!\bigl[F(x_T)-F(x^{\star})\bigr]
\le\frac{L}{2}\Bigl[(1-2\eta\mu E)^{T}\|x_0-x^{\star}\|^{2}
+\frac{\eta\sigma^{2}}{2\mu}\Bigr]. \tag{23}
\]

\textbf{[Step 8]} \emph{Simplify the exponential term.}  
For any $a\in(0,1)$, $a^{T}\le\frac{1}{aT}$. Using $a=1-2\eta\mu E$ and the inequality
$1-a\ge a$ for $a\le\frac12$, we obtain
\[
(1-2\eta\mu E)^{T}\le\frac{1}{2\eta\mu E\,T}. \tag{24}
\]
Insert (24) into (23):
\[
\E\!\bigl[F(x_T)-F(x^{\star})\bigr]
\le\frac{L}{4\eta\mu E\,T}\|x_0-x^{\star}\|^{2}
+\frac{L\eta\sigma^{2}}{4\mu}. \tag{25}
\]

\textbf{[Step 9]} \emph{Final bound matching Theorem statement.}  
Recall that $\mu$ may be zero.  When $\mu=0$ we revert to the standard convex case.
In that regime we keep the recurrence (18) without the contraction term,
which yields the classical bound
\[
\E\!\bigl[F(x_T)-F(x^{\star})\bigr]
\le\frac{\|x_0-x^{\star}\|^{2}}{2\eta T}
+\frac{\eta L\sigma^{2}}{K}.
\]
Combining both regimes gives the unified expression presented in Theorem~\ref{thm:main}. \hfill\qedsymbol
\end{proof}

\section*{Rate Derivation (With Constants)}
From Theorem~\ref{thm:main}, set $\eta=\frac{c}{\sqrt{T}}$ with $c\le\frac{1}{L+\mu}$.
Then
\[
\E\!\bigl[F(x_T)-F(x^{\star})\bigr]
\le
\frac{\|x_0-x^{\star}\|^{2}}{2c}\frac{1}{\sqrt{T}}
+\frac{cL\sigma^{2}}{K}\frac{1}{\sqrt{T}}
+\frac{\mu\|x_0-x^{\star}\|^{2}}{2T}.
\]
Thus the dominant term is $\mathcal{O}(1/\sqrt{T})$ when $\mu=0$ and improves to
$\mathcal{O}(1/T)$ when $\mu>0$ because the last term dominates for large $T$.

\section*{Special Cases}
\begin{enumerate}[label=(\alph*)]
\item \textbf{No proximal term ($\mu=0$).}  
The algorithm reduces to vanilla local SGD.  The bound simplifies to
\[
\E[F(x_T)-F(x^{\star})]\le\frac{\|x_0-x^{\star}\|^{2}}{2\eta T}
+\frac{\eta L\sigma^{2}}{K},
\]
recovering the standard $O(1/\sqrt{T})$ rate for convex smooth problems.

\item \textbf{Single local step ($E=1$).}  
The contraction factor $(1-2\eta\mu E)$ becomes $(1-2\eta\mu)$, matching the analysis
of FedProx with one SGD step per round.

\item \textbf{Large $\mu$ (strong proximal regularization).}  
When $\mu\gg L$, the optimal stepsize is $\eta\approx 1/\mu$, yielding the $O(1/T)$
rate characteristic of strongly convex optimisation even though each $f_k$ is only convex.

\item \textbf{Increasing number of clients $K$.}  
The variance term scales as $\sigma^{2}/K$, showing the usual linear speed‑up with
more participating devices.
\end{enumerate}

\end{document}